\documentclass[UTF8,a4paper,10pt]{ctexart}
\usepackage[a4paper,margin=1.55cm]{geometry}
\usepackage{amsmath,amssymb}
\usepackage{array,longtable,booktabs}
\usepackage{xcolor}
\usepackage{listings}
\usepackage{hyperref}
\hypersetup{colorlinks=true,linkcolor=blue,urlcolor=blue}
\setlength{\parindent}{0pt}
\setlength{\parskip}{4pt}
\setcounter{secnumdepth}{0}
\lstset{basicstyle=\ttfamily\small,backgroundcolor=\color{gray!8},frame=single,breaklines=true,columns=fullflexible}
\begin{document}
\title{GuardFed-AD2+ 完整超参数、攻击空间与三计划实验协议}
\maketitle
\section{1. 给导师的核心结论}
当前实验需要区分三类参数：

\begin{enumerate}
\item \textbf{固定训练协议}：数据集、客户端数量、训练轮数、学习率、batch size、root/server 数据比例等。同一组实验内统一固定，并没有对所有训练参数做笛卡尔积网格。
\item \textbf{实际执行的攻击校准 grid}：F Flip 搜索 5 个候选模式；FedSA 搜索 4 个 gain 与 3 个 norm\_ratio，共 12 组组合。
\item \textbf{AD2+ 内部自适应候选机制}：AD2+ 每轮生成 10 个候选评分视角，用 clean root/server 数据做 one-step 评估，在 accuracy floor 约束下选择一个候选。它不是从外部 baseline 中选择最优，也不是事后根据测试集结果挑选。
\end{enumerate}
准确的论文表述是：\textbf{AD2+ 使用固定的基础训练协议，并在每轮通过 clean root/server 数据对预定义的内部候选族进行自适应选择；攻击强度则通过独立 calibration grid 统一确定。}

\section{2. 最新 attack-strength 三计划的固定训练协议}
{\small
\renewcommand{\arraystretch}{1.18}
\begin{longtable}{|p{0.220\linewidth}|p{0.220\linewidth}|p{0.220\linewidth}|p{0.220\linewidth}|}
\hline
参数 & 当前采用值 & 允许/声明范围 & 作用 \\\hline
\endhead
数据集 & Adult、COMPAS & 固定 & 两个公平联邦学习数据集 \\\hline
客户端数 & 20 & 正整数 & 每轮客户端总数 \\\hline
主实验恶意客户端数 & 4 & 0..20 & 主实验为 20\%  恶意客户端 \\\hline
恶意比例实验 & 2、4、6、8、10 & 10\% ..50\%  & 对应 10\% 、20\% 、30\% 、40\% 、50\%  \\\hline
IID Dirichlet alpha & 5000 & 大于 0 & 近似 IID 客户端划分 \\\hline
non-IID Dirichlet alpha & 5 & 大于 0 & 非 IID 客户端划分 \\\hline
训练轮数 & 70 & 正整数 & 最终指标评估轮 \\\hline
local epochs & 1 & 正整数 & 每个客户端每轮本地训练次数 \\\hline
batch size & 256 & 正整数 & 本地 mini-batch 大小 \\\hline
learning rate & 0.005 & 大于 0 & 本地 Adam 学习率 \\\hline
optimizer & Adam & Adam 或 SGD & 当前核心 runner 使用 Adam \\\hline
clean root/server ratio & 10\%  & 0..1 & 最新攻击强度和三计划实验使用 \\\hline
synthetic root ratio & 0\%  & 0..1 & 最新三计划不使用 synthetic root \\\hline
synthetic method & none & none、gaussian\_copula、bootstrap、noisy\_bootstrap、smote、pca、ctgan 等 & 三计划不生成 synthetic root \\\hline
synthetic epochs & 50 & 正整数 & 仅 synthetic generator 开启时使用 \\\hline
server sampling & stratified\_sensitive & 离散选项 & 按敏感组分层抽取 clean root/server data \\\hline
server alpha & \texttt{None} & optional float & 三计划不额外改变 server Dirichlet alpha \\\hline
server target sensitive & \texttt{None} & optional target & 三计划不强制指定 server sensitive proportion \\\hline
server target label & \texttt{None} & optional target & 三计划不强制指定 server label proportion \\\hline
model input & 不包含 label 和 sensitive feature & 布尔开关 & 所有方法统一，避免算法间输入不一致 \\\hline
aggregation weighting & count & count/equal & 按客户端样本数加权 \\\hline
seeds & 123, 456, 789, 1001, 2024, 3141, 4242, 5050, 6060, 7070 & 整数列表 & 正式结果使用 10 个 seeds \\\hline
\end{longtable}
}
\subsection{关于 5\%  与 10\%  root/server data}
工作区同时保留两条实验轨迹：

\begin{itemize}
\item 旧 paper-table runner 默认 server\_ratio=0.05，对应早期 Table II/III 复现实验；
\item 最新 attack-strength 三计划采用 server\_ratio=0.10、synthetic\_ratio=0.0。
\end{itemize}
这两条协议不能混写在同一张结果表中。新的实验表应在标题或表注中明确写出 10\%  clean root/server。

\section{2.1 三个扩展计划分别使用的设置}
下面这张表是三项扩展计划的实验单元定义。三项计划共享本节第 2 节的训练、数据和 seed 协议；差异只在方法范围、攻击类型和恶意比例设置。

{\scriptsize
\renewcommand{\arraystretch}{1.18}
\begin{longtable}{|p{0.130\linewidth}|p{0.130\linewidth}|p{0.130\linewidth}|p{0.130\linewidth}|p{0.130\linewidth}|p{0.130\linewidth}|p{0.130\linewidth}|}
\hline
计划 & 目的 & 方法范围 & 攻击/对照 & 恶意客户端 & 数据集/分布 & seed 汇报 \\\hline
\endhead
Plan 01 & 强化 F Flip 公平性攻击 & 当前完整方法列表，包含 GuardFed-AD2+ & Benign + F Flip，统一 \texttt{fflip\_mode=all\_unprivileged} & 主实验 4/20 & Adult/COMPAS × IID/non-IID & 10 seeds，保存 raw seed，汇总 mean \\\hline
Plan 02 & 强化 FedSA 性能攻击 & 当前完整方法列表，包含 GuardFed-AD2+ & Benign + FedSA，统一 \texttt{gain=4.5}、\texttt{norm\_ratio=3.0} & 主实验 4/20 & Adult/COMPAS × IID/non-IID & 10 seeds，保存 raw seed，汇总 mean \\\hline
Plan 03 & AD2+ 恶意比例敏感性 & 仅 GuardFed-AD2+ & S-DFA 与 Sp-DFA & 2/4/6/8/10，即 10\% --50\%  & Adult/COMPAS × IID/non-IID & 每个比例 10 seeds，汇总 mean \\\hline
\end{longtable}
}
Plan 01 和 Plan 02 的当前完整方法列表包含 22 个方法：FedAvg、FairFed、Median、FLTrust、FairGuard、FLTrust+FairGuard、GuardFed、FLGMM、FLAURA、LayerGuard、SmartFL、FLTG、FedDNA、LASA、Fed-NGA、Huber-BRFL、LoGoFair、AdaAggRL、FedAMM、FedAA、GuardFed-AD2 和 GuardFed-AD2+。Class-B FL 与 GuardFed-ACT 不纳入统计。

Plan 03 的比例序列固定为：

\begin{lstlisting}
malicious_ratio = {0.10, 0.20, 0.30, 0.40, 0.50}
malicious_clients = {2, 4, 6, 8, 10}
\end{lstlisting}
三个计划的结果表应描述为相同协议下的 10-seed mean。每个 seed 的 ACC、AEOD、ASPD、攻击审计和逐轮轨迹均保留在明细文件中；导师版超参数说明不需要展开某个指标使用哪一轮或哪个 seed。

\section{2.2 模型、随机性和本地优化细节}
5090 runner 实际使用的模型是一个二分类 SimpleMLP：

\begin{lstlisting}
input dimension = dataset feature count
Linear(input_dim, 16)
ReLU
Linear(16, 2)
\end{lstlisting}
其他实现级设置如下：

{\small
\renewcommand{\arraystretch}{1.18}
\begin{longtable}{|p{0.300\linewidth}|p{0.300\linewidth}|p{0.300\linewidth}|}
\hline
参数 & 实际 setting & 说明 \\\hline
\endhead
model & SimpleMLP & 所有方法使用同一模型结构 \\\hline
hidden dimension & 16 & 单个隐藏层 \\\hline
output dimension & 2 & 二分类 \\\hline
loss & weighted cross-entropy & 使用客户端重加权时应用 sample weights \\\hline
client optimizer & Adam & learning rate=0.005 \\\hline
server/root optimizer & Adam & learning rate=0.005 \\\hline
local epochs & 1 & 每个客户端每轮一次本地 epoch \\\hline
seed control & Python/NumPy/PyTorch/CUDA & 每个 seed 统一设置 \\\hline
cuDNN benchmark & \texttt{true} & 5090 runner 实际开启 \\\hline
DataLoader shuffle & \texttt{true} & 客户端和 server/root 本地训练均使用 \\\hline
\end{longtable}
}
\section{2.3 三计划之外的完整攻击场景}
runner 的攻击枚举包括：

\begin{lstlisting}
Benign, F Flip, FOE, S-DFA, Sp-DFA, FedSA
\end{lstlisting}
其中 Plan 01 重点汇总 F Flip，Plan 02 重点汇总 FedSA，Plan 03 汇总 S-DFA/Sp-DFA 的恶意比例变化；FOE 作为原始性能攻击实现和 DFA 性能侧实现保留在 raw result 与 audit 中。

\section{3. 实际运行过的外部 grid}
\subsection{3.1 F Flip 校准 grid：5 个候选模式}
实际候选为：

\begin{lstlisting}
invert
label_conditioned
label_conditioned_reverse
all_privileged
all_unprivileged
\end{lstlisting}
校准 seed 为 314159。结果文件记录的统一选择是：

\begin{lstlisting}
fflip_mode = all_unprivileged
\end{lstlisting}
选择规则是用统一 FedAvg calibration 批次比较公平性影响，并排除无效或退化运行。正式主表使用同一 F Flip 配置，不为不同 baseline 单独改攻击参数。

\subsection{3.2 FedSA 校准 grid：12 个强度组合}
两个维度为：

\begin{lstlisting}
gain       = {1.75, 2.5, 3.5, 4.5}
norm_ratio = {2.0, 3.0, 4.0}
\end{lstlisting}
因此实际运行了 4 x 3 = 12 组组合。最终选择为：

\begin{lstlisting}
fedsa_gain       = 4.5
fedsa_norm_ratio = 3.0
\end{lstlisting}
攻击校准文件还保留了每组候选的 calibration score，可以审计为什么选中这一组。

\subsection{3.3 FOE / 性能侧攻击参数}
5090 代码保留原 Git 实现的 \texttt{attack\_acc\_0.5} 语义，核心缩放常数为：

\begin{lstlisting}
FOE_SCALE = -0.5
\end{lstlisting}
FOE 的实现模式为离散候选：

{\small
\renewcommand{\arraystretch}{1.18}
\begin{longtable}{|p{0.300\linewidth}|p{0.300\linewidth}|p{0.300\linewidth}|}
\hline
\texttt{foe\_mode} & 操作 & 使用情况 \\\hline
\endhead
\texttt{state} & 将恶意客户端 local state 的参数乘以 -0.5 后再形成 update & 默认 FOE mode；当前三计划的基础配置 \\\hline
\texttt{delta} & 将 local update delta 乘以 -0.5 & runner 支持的替代实现 \\\hline
\texttt{zero} & 将恶意 update 置为零 & runner 支持的替代实现 \\\hline
\texttt{fedsa} & 使用 FedSA 的方向偏移和范数约束 & FedSA、S-DFA、Sp-DFA 的性能侧实现 \\\hline
\end{longtable}
}
当前三计划 runner 的明确配置为：

\begin{lstlisting}
foe_mode       = state
sdfa_foe_mode  = fedsa
spdfa_foe_mode = fedsa
\end{lstlisting}
因此，S-DFA 和 Sp-DFA 中的性能侧不是另设一个未记录的攻击强度，而是调用统一的 FedSA 参数：\texttt{gain=4.5}、\texttt{norm\_ratio=3.0}。

\subsection{3.4 F Flip 的数据修改范围}
F Flip 的候选模式是离散集合：

\begin{lstlisting}
{invert, label_conditioned, label_conditioned_reverse,
 all_privileged, all_unprivileged}
\end{lstlisting}
所有 F Flip 模式都满足：

\begin{itemize}
\item 只修改恶意客户端的敏感属性；
\item 不修改真实标签；
\item audit 记录敏感属性实际变化比例；
\item audit 记录 \texttt{label\_changed\_count}，正式配置要求为 0；
\item 正式三计划统一使用 \texttt{all\_unprivileged}。
\end{itemize}
\subsection{3.5 FedSA 的实际更新公式参数}
对恶意客户端原始 update \texttt{g}，代码使用 clean server update 方向进行性能偏移：

\begin{lstlisting}
post_vec = pre_vec - gain × ||pre_vec|| × clean_direction
\end{lstlisting}
随后使用：

\begin{lstlisting}
max_norm = max(||pre_vec||, norm_ratio × ||pre_vec||)
\end{lstlisting}
对超过上限的 update 做范数缩放。所有方法共享同一组 \texttt{gain} 和 \texttt{norm\_ratio}，攻击审计记录 pre/post norm ratio 与 cosine 变化。

\section{4. Baseline-specific settings}
除各方法自身的聚合规则外，所有 baseline 共用同一个数据预处理、模型、optimizer、客户端划分、攻击配置和指标实现。代码中记录的 baseline 相关参数为：

{\small
\renewcommand{\arraystretch}{1.18}
\begin{longtable}{|p{0.300\linewidth}|p{0.300\linewidth}|p{0.300\linewidth}|}
\hline
方法/组件 & 参数 & 实际 setting \\\hline
\endhead
FedAvg & aggregation weighting & \texttt{count} \\\hline
FairFed & \texttt{fairfed\_beta} & 1.0 \\\hline
FairFed & \texttt{use\_reweighting} & \texttt{true} \\\hline
FLTrust & \texttt{trust\_threshold} & 0.2 \\\hline
FairGuard & \texttt{fairguard\_mode} & \texttt{server\_aeod} \\\hline
GuardFed & \texttt{guardfed\_fairness\_lambda} & 20.0 \\\hline
GuardFed & trust selection threshold & 0.2 \\\hline
Median & coordinate aggregation & coordinate-wise median \\\hline
AD2/AD2+ & root calibration & enabled \\\hline
\end{longtable}
}
新增的 FLGMM、FLAURA、LayerGuard、SmartFL、FLTG、FedDNA、LASA、Fed-NGA、Huber-BRFL、LoGoFair、AdaAggRL、FedAMM 和 FedAA 使用仓库中的对应 core reproduction implementation；它们不改变统一训练协议。

\section{5. AD2+ 的实际参数与动态机制}
\subsection{5.1 动态公平风险和 dual multiplier}
对客户端 i，当前轮从 clean root/server 数据得到公平风险 r\_i。公平预算为：

\[B=0.06\].
公平违反项为：

\[v_i=\max(0,r_i-B)\].
当前轮平均风险为 r\_bar，temperature 为 0.35，base weight 为 1.0 时：

\[\lambda_t=1.0\cdot\operatorname{softplus}\left((\bar r-B)/0.35\right)\].
所以 0.06 是固定的公平风险预算，不是每轮自动变化的超参数；但 r\_i、v\_i、r\_bar、lambda\_t 都会随当前模型和当前客户端更新变化。因此 AD2+ 的 adaptive 体现在：公平惩罚、客户端评分、保留集合和聚合权重都会逐轮变化。

\subsection{5.2 AD2 基础参数}
{\small
\renewcommand{\arraystretch}{1.18}
\begin{longtable}{|p{0.300\linewidth}|p{0.300\linewidth}|p{0.300\linewidth}|}
\hline
参数 & 基础值 & AD2+ 内部候选范围 \\\hline
\endhead
act\_fairness\_budget & 0.06 & 固定公平预算 \\\hline
act\_fairness\_metric & aeod\_aspd & aeod、aspd、aeod\_aspd、max \\\hline
act\_temperature & 0.35 & 0.20 或 0.35 \\\hline
act\_keep\_ratio & 0.80 & 0.70、0.80、0.90 \\\hline
act\_risk\_weight & 1.00 & 0.60、0.75、0.80、0.85、0.90、1.10、1.30 \\\hline
act\_violation\_weight & 1.00 & 0.10、0.20、0.25、0.35、0.50 \\\hline
ad2\_utility\_weight & 1.00 & 0.60、0.70、0.80、1.00、1.20 \\\hline
ad2\_centrality\_weight & 0.35 & 0.20、0.25、0.35、0.50 \\\hline
ad2\_alignment\_weight & 0.35 & 0.20、0.25、0.35、0.50 \\\hline
ad2\_norm\_mode & adaptive & adaptive；AD2+ candidate 内部使用 root \\\hline
ad2\_score\_clip & 5.0 & AD2+ candidate 内部为 0 \\\hline
ad2\_norm\_clip\_scale & 2.5 & 固定正数 \\\hline
\end{longtable}
}
AD2+ 顶层 calibration 参数也固定记录如下：

{\small
\renewcommand{\arraystretch}{1.18}
\begin{longtable}{|p{0.300\linewidth}|p{0.300\linewidth}|p{0.300\linewidth}|}
\hline
参数 & 实际 setting & 作用 \\\hline
\endhead
\texttt{ad2\_calibration\_enabled} & \texttt{true} & 开启 clean-root calibration \\\hline
\texttt{ad2\_calibration\_base\_weight} & 1.0 & calibration 的基础权重 \\\hline
\texttt{ad2\_calibration\_budget} & 0.06 & calibration fairness budget \\\hline
\texttt{ad2\_calibration\_temperature} & 0.03 & threshold calibration 温度 \\\hline
\texttt{ad2\_calibration\_quantiles} & 41 & 每个敏感组的 quantile 候选数 \\\hline
\texttt{ad2\_calibration\_max\_acc\_drop} & 0.03 & 配置层面的最大 ACC drop \\\hline
\texttt{ad2\_calibration\_objective} & \texttt{acc\_floor} & 普通 calibration 目标 \\\hline
\texttt{ad2\_plus\_mode} & \texttt{adaptive} & 启用 AD2+ 内部 candidate selector \\\hline
\end{longtable}
}
\subsection{5.3 AD2+ 的 10 个内部候选}
{\tiny
\renewcommand{\arraystretch}{1.18}
\begin{longtable}{|p{0.105\linewidth}|p{0.105\linewidth}|p{0.105\linewidth}|p{0.105\linewidth}|p{0.105\linewidth}|p{0.105\linewidth}|p{0.105\linewidth}|p{0.105\linewidth}|p{0.105\linewidth}|}
\hline
候选 & fairness metric & risk & violation & keep & temp & utility & centrality & alignment \\\hline
\endhead
balanced & aeod\_aspd & 0.75 & 0.20 & 0.80 & 0.35 & 1.00 & 0.35 & 0.35 \\\hline
balanced\_open & aeod\_aspd & 0.75 & 0.20 & 0.90 & 0.35 & 1.00 & 0.35 & 0.35 \\\hline
fair\_stable & aeod\_aspd & 0.90 & 0.25 & 0.80 & 0.35 & 1.00 & 0.35 & 0.35 \\\hline
utility\_fair & aeod\_aspd & 0.60 & 0.10 & 0.80 & 0.35 & 1.20 & 0.50 & 0.50 \\\hline
dual\_strict & aeod\_aspd & 1.10 & 0.35 & 0.80 & 0.35 & 0.70 & 0.25 & 0.25 \\\hline
dual\_sharp & aeod\_aspd & 1.30 & 0.50 & 0.70 & 0.20 & 0.60 & 0.20 & 0.20 \\\hline
aeod\_focus & aeod & 0.85 & 0.25 & 0.80 & 0.35 & 0.80 & 0.35 & 0.35 \\\hline
aspd\_focus & aspd & 0.85 & 0.25 & 0.80 & 0.35 & 0.80 & 0.35 & 0.35 \\\hline
aspd\_strict & aspd & 1.10 & 0.35 & 0.70 & 0.20 & 0.70 & 0.25 & 0.25 \\\hline
max\_guard & max & 0.85 & 0.25 & 0.80 & 0.35 & 0.80 & 0.35 & 0.35 \\\hline
\end{longtable}
}
这 10 行不是 10 个 baseline，也不是事后挑最优结果。每轮都对候选更新做一次 clean-root/server one-step evaluation，在 accuracy floor 约束下选择一个候选。

候选公平损失为：

\[L_{fair}=0.45\,AEOD+0.45\,ASPD+0.10\max(AEOD,ASPD)\].
候选选择分数为：

\[S=ACC-0.35L_{fair}-6.00\max(0,ACC_{floor}-ACC)-0.10\max(0,\max(AEOD,ASPD)-B)\].
其中，只有满足 accuracy floor 的候选优先参与选择；如果没有候选满足，才从全部候选中选择最高分。测试集不参与训练过程中的候选选择。

\section{6. Group threshold calibration 的 grid}
AD2+ 还对两个敏感组分别构造 threshold candidates：

\begin{lstlisting}
quantiles = 41
quantile interval = 0.02 ... 0.98
additional candidates = 0, min(margin)-1e-6, max(margin)+1e-6
\end{lstlisting}
因此每个敏感组最多产生 44 个候选值，去重后进行两组阈值组合搜索；两个组的组合规模最多约为 44 x 44。该搜索使用 clean root/server labels 和 margins，不使用测试集标签。

默认 calibration objective 是 acc\_floor：

\begin{itemize}
\item base accuracy 为 root/server 上未调整阈值的 accuracy；
\item ad2\_calibration\_max\_acc\_drop=0.03 给出最多 3 个百分点的基础上限；
\item AD2+ 每轮进一步把严格候选 floor 收紧到当前候选最高 ACC 下最多 0.5 个百分点的差距；
\item fairness risk 超过 ad2\_calibration\_budget=0.06 时增加 violation 惩罚。
\end{itemize}
\section{7. 哪些参数没有做外部 grid search}
以下参数在当前完整结果中是固定值，而不是已经完成全范围搜索的结果：

\begin{itemize}
\item learning\_rate=0.005；
\item rounds=70；
\item local\_epochs=1；
\item batch\_size=256；
\item server\_ratio=10\% ；
\item act\_fairness\_budget=0.06；
\item ad2\_norm\_clip\_scale=2.5；
\item ad2\_calibration\_max\_acc\_drop=0.03；
\item 10 个正式随机种子集合。
\end{itemize}
这些值有明确的实现范围，但不能写成“我们搜索过所有可能值并选出最优”。如果导师要求完整的 AD2+ 超参数敏感性实验，应另外运行一个标注清楚的新实验，例如：

\begin{lstlisting}
budget ∈ {0.03, 0.06, 0.10}
temperature ∈ {0.20, 0.35, 0.50}
keep_ratio ∈ {0.70, 0.80, 0.90}
centrality/alignment weight ∈ {0.20, 0.35, 0.50}
norm_clip_scale ∈ {1.5, 2.5, 3.5}
\end{lstlisting}
这会形成一个新的 hyperparameter sensitivity 实验，而不是回溯性地把当前固定配置说成已搜索。

\section{8. 建议放入论文 appendix 的英文表述}
\begin{quote}
We use a fixed training protocol with 20 clients, one local epoch, batch size 256, learning rate 0.005, 70 communication rounds, and 10\%  clean root/server data. The non-IID partition uses Dirichlet alpha 5, while IID uses alpha 5000. Attack strengths are calibrated independently using a five-mode F-Flip grid and a 12-point FedSA grid. GuardFed-AD2+ does not select among external methods. Instead, at each communication round it evaluates ten predefined internal scoring candidates on the clean root set, applies an adaptive accuracy floor and fairness-risk penalty, and selects the candidate update with the highest constrained clean-root score. All fixed values and candidate ranges are reported in the supplementary hyperparameter audit.\\
\end{quote}
\section{9. 可核查文件}
5090 上的实际仓库路径为 \texttt{/home/yannan/workspace/GuardFed}。以下是 5090 上直接核对的源文件；Windows workspace 中的同名副本只作为整理和导出材料：

\begin{itemize}
\item \texttt{/home/yannan/workspace/GuardFed/scripts/run\_attack\_strength\_study.py}
\item \texttt{/home/yannan/workspace/GuardFed/scripts/reproduce\_paper\_tables.py}
\item \texttt{/home/yannan/workspace/GuardFed/results/attack\_strength/attack\_config.json}
\item \texttt{/home/yannan/workspace/GuardFed/results/attack\_strength/raw\_results.jsonl}
\end{itemize}
\begin{itemize}
\item \texttt{E:\textbackslash\{\}OneDrive\textbackslash\{\}文档\textbackslash\{\}GuardFed\textbackslash\{\}.codex\_remote\textbackslash\{\}reproduce\_paper\_tables.py}
\item \texttt{E:\textbackslash\{\}OneDrive\textbackslash\{\}文档\textbackslash\{\}GuardFed\textbackslash\{\}.codex\_transfer\textbackslash\{\}run\_attack\_strength\_study.py}
\item \texttt{E:\textbackslash\{\}OneDrive\textbackslash\{\}文档\textbackslash\{\}GuardFed\textbackslash\{\}results\textbackslash\{\}attack\_strength\textbackslash\{\}attack\_config.json}
\item \texttt{E:\textbackslash\{\}OneDrive\textbackslash\{\}文档\textbackslash\{\}GuardFed\textbackslash\{\}results\textbackslash\{\}attack\_strength\textbackslash\{\}raw\_results.jsonl}
\item \texttt{E:\textbackslash\{\}OneDrive\textbackslash\{\}文档\textbackslash\{\}GuardFed\textbackslash\{\}results\textbackslash\{\}attack\_strength\textbackslash\{\}seed\_results.csv}
\item \texttt{E:\textbackslash\{\}OneDrive\textbackslash\{\}文档\textbackslash\{\}GuardFed\textbackslash\{\}results\textbackslash\{\}attack\_strength\textbackslash\{\}audit.csv}
\end{itemize}\end{document}
